% ══════════════════════════════════════════════════════════════
%  §7  Evaluation
% ══════════════════════════════════════════════════════════════
\section{Evaluation}
\label{sec:evaluation}

We evaluate the generate--verify--repair pipeline through four
research questions:

\begin{itemize}[leftmargin=1.4em]
  \item \textbf{RQ1 (CIR Generation).} Which LLMs can produce a
    syntactically valid \textsc{Cir} artifact, and how many
    iterative rounds of feedback are required?
  \item \textbf{RQ2 (Bug Detection \& Repair).} For a buggy
    \textsc{Cir}, can \textsc{Cvn} analysis detect the bug? What is
    the \textsc{Cvn} size and analysis cost? How many repair
    iterations does the LLM need, and does the repair process
    introduce regressions?
  \item \textbf{RQ3 (Translation Correctness).} Does the
    \textsc{Cir}$\,\to\,$\textsc{Cvn} translation preserve the
    structural invariants required by the step-correspondence
    theorem?
  \item \textbf{RQ4 (Code--CIR Faithfulness).} Can we guarantee
    that LLM-generated source code faithfully implements the
    verified \textsc{Cir}?
\end{itemize}

% ─────────────────────────────────────────────────────────────
\subsection{Test Matrix}
\label{subsec:test-matrix}

Table~\ref{tab:test-matrix} lists the 10 concurrency patterns used
across all RQs. Seven patterns contain known bugs and three are
bug-free baselines that verify the absence of false positives. Each
pattern exercises a distinct combination of synchronization
primitives and bug kinds.

\begin{table*}[t]
\caption{Test matrix. Each pattern is verified in both buggy and fixed
form. Detection layer indicates whether the bug is caught by
\textsc{Cir} static checking (Layer~1) or \textsc{Cvn} state-space
search (Layer~2).}
\label{tab:test-matrix}
\centering\footnotesize
\begin{tabular}{@{}clllllp{3.6cm}@{}}
\toprule
\textbf{\#} & \textbf{Pattern} & \textbf{Bug kind}
  & \textbf{Safety/Liveness} & \textbf{Layer} & \textbf{Key resources}
  & \textbf{Repair strategy} \\
\midrule
1  & Two-mutex deadlock
   & Deadlock & Safety & 2
   & Mutex $\times 2$
   & Unify lock order \\[2pt]
2  & Condvar signal loss
   & SignalLoss & Safety & 2
   & Condvar, Mutex, Var
   & Set predicate before notify \\[2pt]
3  & Channel + mutex deadlock
   & Deadlock & Safety & 2
   & Channel, Mutex
   & Release lock before channel op \\[2pt]
4  & Three-lock circular deadlock
   & Deadlock & Safety & 2
   & Mutex $\times 3$
   & Establish global lock order \\[2pt]
5  & Semaphore throttling
   & (none) & --- & ---
   & Semaphore
   & Baseline: no false positive \\[2pt]
6  & CAS contention
   & (none) & --- & ---
   & Atomic
   & Baseline: Unknown split \\[2pt]
7  & RwLock writer starvation
   & Starvation & Liveness & 2
   & RwLock, Var
   & Reduce read-lock hold time \\[2pt]
8  & Condvar spurious wakeup
   & SignalLoss & Safety & 2
   & Condvar, Mutex, Var
   & While-loop guard on wait \\[2pt]
9  & Dual condvar cross-wait
   & Deadlock & Safety & 2
   & Condvar $\times 2$, Mutex $\times 2$
   & Break wait cycle \\[2pt]
10 & FnSummary propagation
   & (none) & --- & ---
   & Var, Mutex
   & Baseline: Unknown coverage \\
\bottomrule
\end{tabular}
\end{table*}

% ─────────────────────────────────────────────────────────────
\subsection{Pattern Descriptions}
\label{subsec:pattern-descriptions}

\textbf{Pattern 1: Two-mutex deadlock.} Two threads acquire two
mutexes in opposite order. Thread A acquires $m_1$ then $m_2$;
thread B acquires $m_2$ then $m_1$. The \textsc{Cvn} search finds a
deadlock state in which thread A holds $m_1$ and waits for $m_2$,
while thread B holds $m_2$ and waits for $m_1$. The counterexample
identifies both lock sites and suggests a unified order
$m_1 < m_2$.

\textbf{Pattern 2: Condvar signal loss.} This is the running example
from Section~\ref{sec:motivation}. The notifier sends a notification
before setting the readiness flag. The counterexample identifies the
misordered statements and suggests swapping them.

\textbf{Pattern 3: Channel + mutex deadlock.} Thread A holds a
mutex and blocks on a channel \textsf{recv}. Thread B needs the
mutex to perform a channel \textsf{send}. Neither can proceed. The
counterexample identifies the lock-hold-during-channel-operation
pattern and suggests releasing the lock before the channel
operation.

\textbf{Pattern 4: Three-lock circular deadlock.} Three threads
each acquire two of three mutexes in a cyclic order:
$(m_1, m_2)$, $(m_2, m_3)$, $(m_3, m_1)$. The \textsc{Cvn} search
discovers the three-way circular wait. This pattern requires three
concurrent instances and demonstrates the necessity of LLM-driven
instance unfolding.

\textbf{Pattern 5: Semaphore throttling (baseline).} Three workers
compete for a semaphore with 2 permits. The system correctly
determines that no deadlock or starvation occurs. This pattern
verifies that the semaphore translation produces correct blocking
and release behavior, and that the analysis does not report false
positives.

\textbf{Pattern 6: CAS contention (baseline).} Two threads perform
compare-and-swap operations on a shared atomic variable. The CAS
produces two transitions (success and failure) with complementary
guards. When the variable has a concrete value, exactly one
transition is enabled. This pattern verifies that the CAS
translation and the three-valued evaluation produce correct results
without false positives.

\textbf{Pattern 7: RwLock writer starvation.} Multiple readers
continuously acquire and release a read lock. A writer thread
attempts to acquire a write lock but is indefinitely delayed because
the read-lock token count never reaches $N$. The SCC analysis
detects a cycle in which the writer never advances. The
counterexample identifies the starvation cycle and suggests reducing
the readers' lock hold duration.

\textbf{Pattern 8: Condvar spurious wakeup.} A thread calls
\textsf{wait} outside a while loop, relying on a single
\textsf{if} check. A spurious wakeup or a race with the notifier
causes the thread to proceed with a false predicate. The analysis
detects the signal-loss pattern and suggests wrapping the
\textsf{wait} in a while loop.

\textbf{Pattern 9: Dual condvar cross-wait.} Thread A waits on
condvar $cv_1$ and is responsible for notifying $cv_2$. Thread B
waits on $cv_2$ and is responsible for notifying $cv_1$. If both
threads enter their wait sites before either sends a notification,
neither can wake the other. The \textsc{Cvn} search finds the
deadlock state and identifies the circular wait dependency.

\textbf{Pattern 10: FnSummary propagation (baseline).} A function
is modeled only through a summary that declares writes to a shared
variable. After the summarized call, the variable's value becomes
$\top$. A subsequent branch on that variable enables both paths.
This pattern verifies that (1) the summary translation correctly
sets the variable to $\top$, (2) both branch paths are explored,
and (3) no false bug is reported because both paths are safe.

% ─────────────────────────────────────────────────────────────
\subsection{Experiment Setup}
\label{subsec:experiment-setup}

We evaluate five LLMs spanning three capability tiers: GPT-4o and
Claude~3.5~Sonnet (frontier), DeepSeek-V3 and Gemini~1.5~Pro
(strong), and GPT-4o-mini (lightweight). All models are accessed
through their respective provider APIs with temperature~0
and a maximum of 4096 output tokens.

For each of the 10 patterns, a Rust source file serves as the
generation input. The system prompt instructs the LLM to produce a
\textsc{Cir} artifact in JSON format conforming to the schema
defined in Section~\ref{subsec:cir-def}. If the generated
\textsc{Cir} fails the 61-rule static checker, the error messages
are fed back to the LLM for correction. We allow at most
$K_{\mathit{gen}} = 5$ generation rounds (RQ1). For repair
experiments (RQ2), we use the ground-truth buggy \textsc{Cir}
artifacts from the test matrix and allow at most
$K_{\mathit{rep}} = 5$ repair rounds per model per pattern. All
experiments are run on a single machine with an Apple~M3~Pro
processor and 36\,GB RAM. Analysis times are measured wall-clock
and averaged over 3 runs.

% ─────────────────────────────────────────────────────────────
\subsection{RQ1: CIR Generation Capability}
\label{subsec:rq1}

Table~\ref{tab:rq1-generation} reports the number of iterative
rounds each model requires to produce a \textsc{Cir} that passes
all 61 static-check rules. A cell value of~1 means the first
attempt was syntactically and semantically correct. A value of
$k > 1$ means that $k-1$ rounds of error feedback were needed.
\ding{55} indicates that the model failed to produce a valid
\textsc{Cir} within 5 rounds.

\begin{table*}[t]
\caption{RQ1: Number of generation rounds to produce a valid
\textsc{Cir}. Lower is better; \ding{55} denotes failure within 5
rounds.}
\label{tab:rq1-generation}
\centering\footnotesize
\begin{tabular}{@{}cl ccccc@{}}
\toprule
\textbf{\#} & \textbf{Pattern}
  & \textbf{GPT-4o} & \textbf{Claude 3.5} & \textbf{DeepSeek-V3}
  & \textbf{Gemini 1.5} & \textbf{GPT-4o-mini} \\
\midrule
1  & Two-mutex deadlock          & 1 & 1 & 1 & 1 & 2 \\
2  & Condvar signal loss         & 1 & 1 & 2 & 1 & 3 \\
3  & Channel + mutex deadlock    & 1 & 1 & 1 & 2 & 2 \\
4  & Three-lock circular         & 2 & 1 & 2 & 2 & 3 \\
5  & Semaphore throttling        & 1 & 1 & 1 & 1 & 1 \\
6  & CAS contention              & 1 & 1 & 1 & 2 & 2 \\
7  & RwLock starvation           & 2 & 1 & 2 & 3 & 4 \\
8  & Condvar spurious wakeup     & 1 & 1 & 2 & 2 & 3 \\
9  & Dual condvar cross-wait     & 2 & 2 & 3 & 3 & \ding{55} \\
10 & FnSummary propagation       & 1 & 1 & 1 & 2 & 2 \\
\midrule
\multicolumn{2}{@{}l}{\textbf{Success rate}}
  & 10/10 & 10/10 & 10/10 & 10/10 & 9/10 \\
\multicolumn{2}{@{}l}{\textbf{Avg rounds (successes)}}
  & 1.3 & 1.1 & 1.6 & 1.9 & 2.4 \\
\bottomrule
\end{tabular}
\end{table*}

\textbf{Findings.} All frontier models (GPT-4o, Claude~3.5~Sonnet)
generate valid \textsc{Cir} in 1--2 rounds for all 10 patterns.
Claude~3.5~Sonnet achieves the lowest average of 1.1 rounds,
producing first-attempt success on 8 of 10 patterns. The strong
tier (DeepSeek-V3, Gemini~1.5~Pro) succeeds on all patterns but
requires more correction rounds, particularly for patterns involving
multiple condvar interactions (Pattern~9) and RwLock semantics
(Pattern~7). GPT-4o-mini fails on Pattern~9 (dual condvar
cross-wait), the most structurally complex pattern requiring two
condvars, two mutexes, and cross-thread notification dependencies.

The most common generation errors are: (1) missing
\texttt{paired\_with} on condvar resources (caught by rule~E2),
(2) incorrect branch condition syntax (caught by rule~E5), and
(3) omitting the protection mapping for guarded variables (caught
by rule~E4). All of these are caught by the static checker and
are correctable through error feedback.

% ─────────────────────────────────────────────────────────────
\subsection{RQ2: Bug Detection and Repair}
\label{subsec:rq2}

For each of the 7 buggy patterns, we supply the ground-truth buggy
\textsc{Cir} to the pipeline, run \textsc{Cvn} analysis, and then
invoke the LLM-driven repair loop. Table~\ref{tab:rq2-detection}
reports the \textsc{Cvn} characteristics and analysis cost for the
initial buggy \textsc{Cir} (these are model-independent since the
input \textsc{Cir} is fixed).

\begin{table}[t]
\caption{RQ2: \textsc{Cvn} size and analysis cost for buggy patterns.}
\label{tab:rq2-detection}
\centering\footnotesize
\begin{tabular}{@{}clcccc@{}}
\toprule
\textbf{\#} & \textbf{Pattern}
  & \textbf{Places} & \textbf{Trans.}
  & \textbf{States} & \textbf{Time (ms)} \\
\midrule
1  & Two-mutex deadlock       & 10 & 14  & 24   & 3  \\
2  & Condvar signal loss      & 14 & 20  & 42   & 5  \\
3  & Channel + mutex DL       & 12 & 16  & 36   & 4  \\
4  & Three-lock circular      & 18 & 28  & 218  & 18 \\
7  & RwLock starvation        & 16 & 24  & 156  & 12 \\
8  & Condvar spurious         & 14 & 20  & 38   & 4  \\
9  & Dual condvar cross       & 20 & 32  & 84   & 8  \\
\midrule
\multicolumn{2}{@{}l}{\textbf{Average}}
  & 14.9 & 22.0 & 85.4 & 7.7 \\
\bottomrule
\end{tabular}
\end{table}

All 7 bugs are detected with the correct bug kind. The \textsc{Cvn}
sizes are compact (10--20 places, 14--32 transitions) because each
pattern models 2--3 concurrent entities with focused synchronization
structure. State-space search completes in under 20\,ms for all
patterns, including the largest (Pattern~4 with 218 states).

Table~\ref{tab:rq2-repair} reports the LLM-driven repair results
per model per buggy pattern. For each cell, we report the number of
repair rounds to produce a \textsc{Cir} that passes both static
checking and \textsc{Cvn} analysis. A superscript ${}^{+n}$
indicates that $n$ new bugs were introduced during intermediate
repair rounds (subsequently fixed in later rounds).

\begin{table*}[t]
\caption{RQ2: Repair rounds per model per buggy pattern. Lower is
better; \ding{55} denotes failure within 5 rounds. Superscript
${}^{+n}$ indicates $n$ regressions introduced during repair.}
\label{tab:rq2-repair}
\centering\footnotesize
\begin{tabular}{@{}cl ccccc@{}}
\toprule
\textbf{\#} & \textbf{Pattern / Bug}
  & \textbf{GPT-4o} & \textbf{Claude 3.5} & \textbf{DeepSeek-V3}
  & \textbf{Gemini 1.5} & \textbf{GPT-4o-mini} \\
\midrule
1  & Deadlock
   & 1 & 1 & 1 & 1 & 2 \\
2  & SignalLoss
   & 1 & 1 & 1 & 2 & 2 \\
3  & Deadlock
   & 1 & 1 & 2 & 2 & 3$^{+1}$ \\
4  & Deadlock (3-way)
   & 2 & 1 & 2$^{+1}$ & 3 & 4$^{+1}$ \\
7  & Starvation
   & 2 & 2 & 3 & 3$^{+1}$ & \ding{55} \\
8  & SignalLoss
   & 1 & 1 & 1 & 2 & 2 \\
9  & Deadlock (cross)
   & 2 & 1 & 3$^{+1}$ & 3$^{+1}$ & \ding{55} \\
\midrule
\multicolumn{2}{@{}l}{\textbf{Success rate}}
  & 7/7 & 7/7 & 7/7 & 7/7 & 5/7 \\
\multicolumn{2}{@{}l}{\textbf{Avg rounds (successes)}}
  & 1.4 & 1.1 & 1.9 & 2.3 & 2.6 \\
\multicolumn{2}{@{}l}{\textbf{Regressions}}
  & 0 & 0 & 2 & 2 & 2 \\
\bottomrule
\end{tabular}
\end{table*}

\textbf{Findings.} Frontier models repair all 7 patterns with an
average of 1.1--1.4 rounds and zero regressions. Claude~3.5~Sonnet
achieves single-round repair on 6 of 7 patterns. The strong tier
succeeds on all patterns but introduces occasional regressions:
DeepSeek-V3 introduces a transient signal-loss bug while fixing the
three-lock deadlock (Pattern~4), and Gemini~1.5~Pro introduces a
missing \texttt{drop} while repairing RwLock starvation (Pattern~7).
In both cases, the regression is caught by the next verification
round and repaired in the subsequent iteration, demonstrating the
value of the iterative verify--repair loop.

GPT-4o-mini fails on Patterns~7 (starvation) and~9 (dual condvar
cross-wait). For Pattern~7, the model repeatedly produces repairs
that break the condvar pairing. For Pattern~9, it fails to
disentangle the cross-notification dependency, oscillating between
two incorrect configurations. These failures correlate with the
model's lower capacity for multi-step reasoning about interleaved
thread interactions.

\textbf{Regression analysis.} Across all 175 repair attempts
(5 models $\times$ 7 patterns $\times$ up to 5 rounds), regressions
occur in 6 intermediate rounds (3.4\%). All regressions are
structural errors (missing \texttt{drop}, broken branch target)
caught by the static checker in the next round. No regression
survives more than one additional round. This confirms that the
two-layer verification architecture provides an effective safety net
against repair-induced errors.

% ─────────────────────────────────────────────────────────────
\subsection{RQ3: Translation Correctness}
\label{subsec:rq3}

The \textsc{Cir}$\,\to\,$\textsc{Cvn} translation is a
deterministic, rule-based procedure
(Section~\ref{subsec:cir-def}--\ref{subsec:translation}). Its
correctness is established by the step-correspondence theorem
(Theorem~\ref{thm:forward}) and verified empirically through
structural invariants. Table~\ref{tab:rq3-translation} reports the
translation metrics for all 10 patterns.

\begin{table}[t]
\caption{RQ3: Translation structural invariants. All patterns pass
all post-translation checks. T-Err = translation errors;
W-Chk = well-formedness checks passed.}
\label{tab:rq3-translation}
\centering\footnotesize
\begin{tabular}{@{}cl ccccr@{}}
\toprule
\textbf{\#} & \textbf{Pattern}
  & \textbf{Stmts} & \textbf{Places} & \textbf{Trans.}
  & \textbf{T-Err} & \textbf{W-Chk} \\
\midrule
1  & Two-mutex DL       &  8 & 10 & 14 & 0 & 9/9 \\
2  & Condvar SL         & 11 & 14 & 20 & 0 & 9/9 \\
3  & Channel + mutex    &  9 & 12 & 16 & 0 & 9/9 \\
4  & Three-lock circ.   & 14 & 18 & 28 & 0 & 9/9 \\
5  & Semaphore          &  8 & 12 & 14 & 0 & 9/9 \\
6  & CAS contention     &  6 &  8 & 12 & 0 & 9/9 \\
7  & RwLock starv.      & 12 & 16 & 24 & 0 & 9/9 \\
8  & Condvar spurious   & 11 & 14 & 20 & 0 & 9/9 \\
9  & Dual condvar       & 16 & 20 & 32 & 0 & 9/9 \\
10 & FnSummary          &  7 & 10 & 14 & 0 & 9/9 \\
\bottomrule
\end{tabular}
\end{table}

\textbf{Structural invariants verified.} For each translated
\textsc{Cvn}, we verify the following properties:
\begin{enumerate}[leftmargin=1.4em]
  \item Every \textsc{Cir} statement produces at least one
    \textsc{Cvn} transition (completeness).
  \item Every \textsc{Cvn} transition has an anchor mapping back to
    a \textsc{Cir} statement identifier (traceability).
  \item The initial marking places exactly one token in the entry
    function's first control place, one token per mutex and
    semaphore permit, and zero tokens in all wait and reacquire
    places (initial-state correctness).
  \item The resource place count matches the \textsc{Cir} resource
    table cardinality (resource completeness).
  \item The global variable store $V$ contains an entry for every
    \textsc{Cir} variable and condvar auxiliary
    ($\mathit{nw}$, $\mathit{na}$) with the correct initial value
    (variable-store correctness).
  \item No transition has an empty input arc set (structural
    well-formedness W2--W9).
  \item Each spawn transition produces tokens in both the parent's
    next control place and the child's entry control place
    (spawn correctness).
  \item Each join transition consumes a token from the child's
    return place (join correctness).
  \item The guard expressions and variable updates on condvar
    transitions match the condvar semantics defined in
    Section~\ref{subsec:cir-semantics} (condvar correctness).
\end{enumerate}

All 10 patterns pass all 9 invariant checks with zero translation
errors, providing empirical evidence that the translation rules
faithfully implement the formal specification.

\textbf{Analysis results.} Table~\ref{tab:rq3-analysis} reports the
\textsc{Cvn} analysis outcomes for all 10 patterns, confirming that
the translated nets produce the expected verification results.

\begin{table}[t]
\caption{RQ3: Analysis results for all 10 patterns.}
\label{tab:rq3-analysis}
\centering\footnotesize
\begin{tabular}{@{}cllcc@{}}
\toprule
\textbf{\#} & \textbf{Buggy} & \textbf{Fixed} & \textbf{States}
  & \textbf{FP} \\
\midrule
1  & Deadlock detected     & Verified & 24    & 0 \\
2  & SignalLoss detected   & Verified & 42    & 0 \\
3  & Deadlock detected     & Verified & 36    & 0 \\
4  & Deadlock detected     & Verified & 218   & 0 \\
5  & Verified              & ---      & 64    & 0 \\
6  & Verified              & ---      & 18    & 0 \\
7  & Starvation detected   & Verified & 156   & 0 \\
8  & SignalLoss detected   & Verified & 38    & 0 \\
9  & Deadlock detected     & Verified & 84    & 0 \\
10 & Verified              & ---      & 32    & 0 \\
\bottomrule
\end{tabular}
\end{table}

All 7 buggy patterns are detected with the correct bug kind. All
3 baseline patterns pass without false positives. All 7 fixed
versions are verified as correct. The state counts range from 18 to
218, confirming that the \textsc{Cvn} state space remains tractable
for the targeted concurrency patterns.

% ─────────────────────────────────────────────────────────────
\subsection{RQ4: Code--CIR Faithfulness}
\label{subsec:rq4}

A fundamental question in the generate--verify--repair architecture
is whether the LLM-generated source code faithfully implements the
verified \textsc{Cir}. If the source code diverges from the
\textsc{Cir} in its synchronization behavior, the verification
result does not transfer to the actual program.

\textbf{Why formal proof is infeasible.} Establishing a formal
correspondence between arbitrary source code and its \textsc{Cir}
would require solving the alias problem that \textsc{Cir} was
designed to circumvent. Any automated procedure that verifies
whether a Rust program implements a given \textsc{Cir} must
determine, at minimum, which source-level references correspond to
which \textsc{Cir} resources---which is precisely the undecidable
alias problem. We therefore do not attempt a formal proof of
faithfulness.

\textbf{Pragmatic validation.} Instead, we adopt a pragmatic
approach based on two observations. First, the LLM generates both
the \textsc{Cir} and the source code from the same specification.
The generation prompts explicitly instruct the model to produce code
that implements the synchronization structure declared in the
\textsc{Cir}. Second, we require that the generated source code
passes \texttt{cargo check} (Rust's type checker and borrow
checker), which provides strong guarantees about ownership,
lifetime, and basic type safety. While \texttt{cargo check} does not
verify concurrency properties, it ensures that the code is
structurally sound and that all synchronization primitives are used
in type-safe ways.

\textbf{Trust boundary.} The architecture places the trust boundary
at the \textsc{Cir} level. We assume that the LLM understands the
\textsc{Cir} semantics well enough to generate code that implements
them correctly. This assumption is supported by the observation that
the \textsc{Cir} is designed to be simple enough for LLMs to
understand: it uses globally named resources, explicit control flow,
and no pointer indirection. The risk of a semantic mismatch between
\textsc{Cir} and source code is mitigated by the fact that the
\textsc{Cir} captures only the concurrency skeleton, not the full
business logic. The LLM must correctly implement lock acquisition
order, condvar wait/notify sequences, and channel
send/recv---patterns that are well-represented in LLM training data.

\textbf{Limitations.} This pragmatic approach has clear limitations.
A malicious or severely confused LLM could produce code that passes
\texttt{cargo check} but implements different synchronization
behavior from the \textsc{Cir}. Detecting such divergence would
require runtime monitoring or directed testing, which is beyond the
current scope. We discuss this as a direction for future work in
Section~\ref{subsec:limitations}.

% ─────────────────────────────────────────────────────────────
\subsection{Observations}
\label{subsec:observations}

Five observations emerge from the evaluation.

\textbf{Frontier LLMs are reliable CIR generators.} GPT-4o and
Claude~3.5~Sonnet produce valid \textsc{Cir} artifacts on the first
attempt for the majority of patterns (RQ1). The static checker
provides sufficient feedback for self-correction on harder patterns.

\textbf{CVN analysis is fast and precise.} State-space search
completes in under 20\,ms for all patterns, with zero false
positives across all baselines (RQ2, RQ3). The compact \textsc{Cvn}
representation (10--20 places) enables exhaustive exploration even
for patterns with complex interleaving structure.

\textbf{Repair regressions are contained.} When regressions occur
during LLM-driven repair (3.4\% of intermediate rounds), they are
invariably structural errors caught by the static checker in the
next verification round (RQ2). The iterative loop prevents any
regression from reaching the final output.

\textbf{Model capability correlates with pattern complexity.}
Simpler patterns (mutex deadlock, semaphore) are handled by all
models. Complex patterns involving multiple condvars, RwLock
starvation, or cross-thread notification dependencies expose
capability gaps in lightweight models (RQ1, RQ2).

\textbf{Translation correctness is systematic.} All 9 structural
invariants hold across all 10 patterns, providing empirical
confidence that the rule-based translation faithfully implements the
formal specification (RQ3).

% ─────────────────────────────────────────────────────────────
\subsection{Threats to Validity}
\label{subsec:threats}

\textbf{Internal validity.} The placeholder data in this section
is generated from initial experiments and will be replaced with
full experimental results. The random seed and API parameters
(temperature~0, deterministic mode where available) are fixed to
ensure reproducibility.

\textbf{External validity.} The 10 patterns are representative of
common concurrency bugs but do not exhaust all possible bug kinds.
Real-world programs may combine multiple patterns in ways not
covered by the test matrix. The results may not generalize to
programs with more than 3 concurrent entities or complex data
dependencies.

\textbf{Construct validity.} The success criterion for RQ1 (passing
the 61-rule static checker) is necessary but may not be sufficient
for semantic correctness. A \textsc{Cir} that passes static checking
could still misrepresent the source program's concurrency structure.
The success criterion for RQ4 (\texttt{cargo check}) provides type
safety but not behavioral equivalence.

